% !TeX root = ../main.tex
\section{Цель работы}

\subsection {Методика измерения коэффициента передачи и развязки с учётом ограничений приборов}

Ограничения оборудования:

Диапазон выходной мощности генератора: $0 \dots +28\text{ дБм}$.

Предельно допустимая мощность на входе анализатора спектра: не более $1\text{ Вт} = +30\text{ дБм}$ (для обеспечения линейного режима смесителя анализатора уровень не должен превышать $-10 \dots 0\text{ дБм}$).

Измерение коэффициента усиления ($S_{21}$):

Генератор настраивается на минимальную мощность $P_{\text{ген}} = 0\text{ дБм}$.

Для предотвращения перегрузки входа усилителя и входа анализатора между генератором и усилителем включается постоянный аттенюатор (например, $40\text{ дБ}$), снижающий мощность на входе усилителя до $-40\text{ дБм}$ (глубокий режим малого сигнала).

Сигнал проходит через кабель 1, усилитель, кабель 2 на анализатор спектра; фиксируется мощность $P_1(f)$.

Калибровка тракта: усилитель исключается из тракта, кабели соединяются переходником $SMA(f)-SMA(f)$, измеряется мощность калибровки $P_2(f)$.

Истинный коэффициент передачи усилителя:


$$K_p(f)\text{ [дБ]} = S_{21}(f)\text{ [дБ]} = P_1(f) - P_2(f)$$


Такая разность исключает нелинейную частотную зависимость затухания кабелей 1 и 2, а также погрешность аттенюатора.

Измерение развязки ($S_{12}$):

Усилитель переворачивается в тракте: сигнал от генератора подаётся на выход усилителя, а ослабленный сигнал снимается с его входа.

Развязка:


$$\text{Reverse Isolation [дБ]} = S_{12}(f)\text{ [дБ]} = P_{\text{вых.обр}}(f) - P_2(f)$$

\subsection {Расчёт мощности шума на выходе при заданном входном шуме}

Дано:

Мощность входного шума: $P_{\text{ш.вх}} = 0\text{ дБм} = 1\text{ мВт} = 10^{-3}\text{ Вт}$;

Полоса шума: $\Pi = 100\text{ МГц} = 10^8\text{ Гц}$;

Коэффициент усиления усилителя: $K_p = +10\text{ дБ} \implies K_p = 10$;

Коэффициент шума усилителя: $K_{\text{ш}} = 3\text{ дБ} \implies K_{\text{ш}} = 10^{3/10} \approx 2{,}0$.

Физический анализ:

Собственная спектральная плотность теплового шума согласованного источника при комнатной температуре: $G_0 = k T_0 = 4{,}0 \cdot 10^{-21}\text{ Вт/Гц} = -174\text{ дБм/Гц}$.

Тепловой шум в полосе $100\text{ МГц}$:


$$P_{\text{тепловой}} = G_0 \Pi = -174 + 10\lg(10^8) = -174 + 80 = -94\text{ дБм}$$

Мощность поданного шума $P_{\text{ш.вх}} = 0\text{ дБм} \gg -94\text{ дБм}$ (на 94 дБ превышает тепловой порог). Внутренний собственный шум усилителя составляет:


$$P_{\text{ш.собств.вых}} = (K_{\text{ш}} - 1) k T_0 \Pi K_p \approx (2 - 1) \cdot 10^{-94/10} \cdot 10 = 10^{-93}\text{ мВт} \approx -93\text{ дБм}$$

Полная выходная мощность шума:


$$P_{\text{ш.вых}} = P_{\text{ш.вх}} \cdot K_p + P_{\text{ш.собств.вых}} \approx P_{\text{ш.вх}} \cdot K_p$$


В децибелах:


$$P_{\text{ш.вых}}\text{ [дБм]} = P_{\text{ш.вх}}\text{ [дБм]} + K_p\text{ [дБ]} = 0\text{ дБм} + 10\text{ дБ} = \mathbf{+10\text{ дБм}}\text{ (10 мВт)}$$

\subsection {Расчёт отношения сигнал/шум на выходе усилителя}

Дано:

Мощность сигнала на входе: $P_{c.\text{вх}} = -80\text{ дБм}$;

Эффективная полоса: $\Pi = 10\text{ МГц} = 10^7\text{ Гц}$;

Коэффициент шума: $K_{\text{ш}} = NF = 6\text{ дБ}$.

Расчёт:

Мощность теплового шума на входе при стандартных условиях ($T_0 = 293\text{ К}$):


$$P_{\text{ш.вх}}\text{ [дБм]} = -174\text{ дБм/Гц} + 10\lg(\Pi) = -174 + 10\lg(10^7) = -174 + 70 = -104\text{ дБм}$$

Входное отношение сигнал/шум:


$$SNR_{\text{вх}}\text{ [дБ]} = P_{c.\text{вх}}\text{ [дБм]} - P_{\text{ш.вх}}\text{ [дБм]} = -80 - (-104) = 24\text{ дБ}$$

По определению коэффициента шума:


$$NF\text{ [дБ]} = SNR_{\text{вх}}\text{ [дБ]} - SNR_{\text{вых}}\text{ [дБ]}$$


Следовательно:


$$SNR_{\text{вых}}\text{ [дБ]} = SNR_{\text{вх}}\text{ [дБ]} - NF\text{ [дБ]} = 24\text{ дБ} - 6\text{ дБ} = \mathbf{18\text{ дБ}}$$

\subsection {Методика измерения коэффициента шума усилителя}

Теоретическое определение:


$$K_{\text{ш}} = \frac{G_{\text{вых.реал}}}{G_0 K_p} = \frac{P_{\text{ш.вых}}}{k T_0 \Pi_{\text{ш}} K_p}$$

Методика косвенного измерения (по схеме с согласованной нагрузкой $50\text{ Ом}$):

Вход усилителя глушится согласованной экранированной нагрузкой $50\text{ Ом}$ при комнатной температуре.

Выход усилителя через кабель с потерями $IL_2$ подключается к анализатору спектра.

На анализаторе спектра выбирается полоса разрешения $RBW$, видеополоса усреднения $VBW \ll RBW$ и детектор среднеквадратичного значения (RMS).

Измеряется мощность шума $P_{\text{изм}}$ [дБм] в полосе $RBW$.

Расчётная формула для логарифмических величин:


$$K_{\text{ш}}\text{ [дБ]} = P_{\text{изм}}\text{ [дБм]} + IL_2\text{ [дБ]} - \left( -174\text{ дБм/Гц} + 10\lg(RBW\text{ [Гц]}) + K_p\text{ [дБ]} \right)$$

\subsection {Настройки анализатора спектра для снижения его собственных шумов}

Собственный шум анализатора спектра (DANL — Displayed Average Noise Level) минимизируется следующими регулировками:

Уменьшение входного аттенюатора анализатора спектра (RF Attenuation): установить $0\text{ дБ}$ (или минимально возможное значение). Каждый децибел входной аттенюации напрямую увеличивает коэффициент шума анализатора спектра на 1 дБ.

Включение встроенного предусилителя (Preamplifier ON): снижает общий коэффициент шума прибора на $15 \dots 25\text{ дБ}$.

Уменьшение полосы пропускания ПЧ ($RBW$): мощность шума пропорциональна $RBW$ ($P_{\text{ш}} \sim 10\lg(RBW)$). Уменьшение $RBW$ в 10 раз опускает шумовую дорожку на $10\text{ дБ}$.

Снижение полосы видеофильтра ($VBW$) и усреднение (Video Averaging): сглаживает шумовые пики и уменьшает дисперсию шумовой дорожки.

\subsection {Методика измерения точки децибельной компрессии $P_{1dB}$}

Определение: $P_{1dB}$ — уровень мощности, при котором коэффициент усиления усилителя уменьшается на $1\text{ дБ}$ по сравнению с линейным (малосигнальным) усилением $K_{p0}$.

Методика:

На фиксированной частоте подаётся входной сигнал от генератора с переменной мощностью (с шагом $1\text{ дБ}$).

Для защиты входа анализатора спектра на выход усилителя ставится калиброванный аттенюатор (например, $10\text{ дБ}$) с известным затуханием $IL_2$.

Фиксируется зависимость выходной мощности $P_{\text{вых}}$ от входной мощности $P_{\text{вх}}$.

Строится линейная экстраполяция: $P_{\text{вых.лин}} = P_{\text{вх}} + K_{p0}$.

Точка, где реальная кривая отстаёт от идеальной прямой на $1\text{ дБ}$, даёт:


$$OP_{1dB} = P_{\text{вых.реал}} = P_{\text{вых.лин}} - 1\text{ дБ}$$

$$IP_{1dB} = OP_{1dB} - (K_{p0} - 1)$$

\subsection {Качественное изменение точки компрессии $P_{1dB}$ с ростом частоты}

С ростом частоты точка децибельной компрессии $P_{1dB}$ широкополосного усилителя монотонно снижается (смещается влево и вниз).

Причины: падение граничной крутизны транзисторов, влияние барьерных емкостей переходов, паразитных индуктивностей выводов корпуса, а также спад коэффициента усиления каскадов, требующий большего размаха входного напряжения для достижения предельного режима по току/напряжению.

\subsection {Методика измерения $IP_3$ и формирование двухтонового сигнала}

Методика: На вход подаются два синусоидальных сигнала равной мощности $P_C$ с близкими частотами $f_1$ и $f_2$ ($\Delta f = 1\text{ МГц}$). Измеряется мощность основного тона на выходе $P_{f_{1,2}}$ и продуктов третьего порядка на частотах $2f_1 - f_2$ и $2f_2 - f_1$ ($IM_3$).


$$IMD_3\text{ [дБ]} = P_{f_{1,2}}\text{ [дБм]} - IM_3\text{ [дБм]}$$

$$OIP_3\text{ [дБм]} = P_{f_{1,2}}\text{ [дБм]} + \frac{IMD_3\text{ [дБ]}}{2}$$

$$IIP_3\text{ [дБм]} = OIP_3\text{ [дБм]} - K_p\text{ [дБ]}$$

Формирование через смеситель: Двухтоновый сигнал можно получить модуляцией несущей $f_0$ синусоидальным сигналом $\Delta f / 2$ в балансном смесителе с подавленной несущей (DSB-SC). На выходе образуются боковые частоты $f_0 \pm \Delta f / 2$.

Почему сумматор мощности предпочтительнее смесителя:

Смеситель сам является существенно нелинейным узлом и создаёт собственный веер комбинационных частот высших порядков, маскирующих искажения исследуемого усилителя.

Использование двух независимых генераторов и пассивного сумматора/делителя мощности обеспечивает высокую развязку между источниками (особенно при установке аттенюаторов по входам сумматора), предотвращая взаимную интермодуляцию в выходных каскадах генераторов.

\subsection {Подбор частот и спектр продуктов третьего порядка ($\Delta f = 1\text{ МГц}$)}

Выбор частот тонов:

$$f_1 = 1000\text{ МГц},\quad f_2 = 1001\text{ МГц}\quad (\Delta f = f_2 - f_1 = 1\text{ МГц})$$

Частоты продуктов интермодуляции третьего порядка ($IM_3$):

$$f_{IM3\_1} = 2f_1 - f_2 = 2 \cdot 1000 - 1001 = \mathbf{999\text{ МГц}}$$

$$f_{IM3\_2} = 2f_2 - f_1 = 2 \cdot 1001 - 1000 = \mathbf{1002\text{ МГц}}$$

Продукты располагаются симметрично по бокам от основных тонов с тем же интервалом $\Delta f = 1\text{ МГц}$.

\subsection {Качественное изменение $IP_3$ с ростом частоты}

С увеличением рабочей частоты точка $IP_3$ широкополосного усилителя снижается. Это обусловлено спадом петлевого усиления отрицательной обратной связи (ООС) в активных элементах, ростом фазовых сдвигов и проявлением нелинейностей паразитных вольт-фарадных характеристик внутренних емкостей переходов транзисторов.

\subsection {Влияние входного аттенюатора $6\text{ дБ}$ на $IIP_3$ и $OIP_3$}

Пусть исходные параметры усилителя: коэффициент усиления $K_p$, входная точка $IIP_{3,0}$, выходная $OIP_{3,0} = IIP_{3,0} + K_p$.

К входу подключён аттенюатор с затуханием $A = 6\text{ дБ}$.

Входная точка компрессии/интермодуляции всей цепочки ($IIP_3$):

Чтобы на вход самого усилителя пришла прежняя мощность, на вход аттенюатора нужно подать сигнал на $6\text{ дБ}$ мощнее.

$$IIP_{3,\text{нов}} = IIP_{3,0} + A = \mathbf{IIP_{3,0} + 6\text{ дБ}}\quad \text{(увеличивается на 6 дБ)}$$

Выходная точка ($OIP_3$):

Новый коэффициент усиления каскада: $K_{p,\text{нов}} = K_p - 6\text{ дБ}$.

$$OIP_{3,\text{нов}} = IIP_{3,\text{нов}} + K_{p,\text{нов}} = (IIP_{3,0} + 6) + (K_p - 6) = IIP_{3,0} + K_p = \mathbf{OIP_{3,0}}\quad \text{(не изменяется)}$$

\subsection {Полоса продуктов интермодуляции третьего порядка для сигнала с полосой $10\text{ МГц}$}

Анализ: Продукт третьего порядка во временной области описывается выражением $y(t) \sim x^3(t)$. В частотной области операция возведения в третью степень эквивалентна трёхкратной свёртке спектров:

$$S_3(f) \sim S(f) * S(f) * S(f)$$

При свёртке двух сигналов с полосами $\Delta F_1$ и $\Delta F_2$ результирующая полоса равна сумме их полос: $\Delta F_{\Sigma} = \Delta F_1 + \Delta F_2$.

Для трёхкратной свёртки полосы суммируются:

$$\Delta F_{IM3} = 3 \cdot \Delta F_{\text{сигн}} = 3 \cdot 10\text{ МГц} = \mathbf{30\text{ МГц}}$$

(Центральная частота остаётся $1\text{ ГГц}$).

\subsection {Сравнение коэффициента шума схемы: аттенюатор на входе vs на выходе}

Дано:

Усилитель: $K_{p1} = +20\text{ дБ} = 100$, $K_{\text{ш1}} = 3\text{ дБ} \approx 2{,}0$ ($1{,}995$);

Аттенюатор: затухание $A = 6\text{ дБ} \implies K_{p2} = -6\text{ дБ} = 0{,}25$, коэффициент шума пассивного аттенюатора равен его ослаблению: $K_{\text{ш2}} = 6\text{ дБ} \approx 4{,}0$ ($3{,}981$).

Случай А: Аттенюатор на входе, усилитель на выходе:

По формуле Фрииса:

$$K_{\text{ш},\Sigma A} = K_{\text{ш.атт}} + \frac{K_{\text{ш.ус}} - 1}{K_{p.\text{атт}}} = 4{,}0 + \frac{2{,}0 - 1}{0{,}25} = 4{,}0 + 4{,}0 = \mathbf{8{,}0}\text{ или в дБ: } 10\lg 8 \approx \mathbf{9{,}03\text{ дБ}})$$

(Аттенюатор на входе напрямую ухудшает $K_{\text{ш}}$ тракта на величину своего затухания).

Случай Б: Аттенюатор переподключён на выход усилителя:

По формуле Фрииса:

$$K_{\text{ш},\Sigma B} = K_{\text{ш.ус}} + \frac{K_{\text{ш.атт}} - 1}{K_{p.\text{ус}}} = 2{,}0 + \frac{4{,}0 - 1}{100} = 2{,}0 + 0{,}03 = \mathbf{2{,}03}\text{ или в дБ: } 10\lg 2{,}03 \approx \mathbf{3{,}07\text{ дБ}})$$

Изменение коэффициента шума:

$$\Delta K_{\text{ш}} = K_{\text{ш},\Sigma A}\text{ [дБ]} - K_{\text{ш},\Sigma B}\text{ [дБ]} \approx 9{,}03\text{ дБ} - 3{,}07\text{ дБ} \approx \mathbf{5{,}96\text{ дБ}} \approx \mathbf{6\text{ дБ}}$$

Вывод: При переносе аттенюатора со входа на выход результирующий коэффициент шума схемы снижается практически точно на величину затухания аттенюатора ($6\text{ дБ}$), кардинально улучшая чувствительность тракта.